% Clase 26 --- Fibra optica y onda evanescente (§6.2 y §6.3).
% (a) la aplicacion: la luz "dobla" a fuerza de reflexiones totales sucesivas.
% (b) la letra chica: del otro lado SI penetra algo, sin propagarse; si el
%     segundo medio es fino, la onda continua (reflexion total frustrada).
\begin{center}
\begin{tikzpicture}[scale=0.95, transform shape]

% ---------------- (a) fibra óptica ----------------
\begin{scope}
  \fill[figblue!10] (0,0) rectangle (6.0,1.2);
  \draw[figline, line width=1pt] (0,0) -- (6.0,0);
  \draw[figline, line width=1pt] (0,1.2) -- (6.0,1.2);
  \node[figlblaux, anchor=west] at (0.05,1.5) {revestimiento: $n_2 < n_1$};

  \draw[figvecres] (-0.9,-0.36) -- (0,0.2);
  \draw[figred, line width=1.1pt]
    (0,0.2) -- (1.25,1.2) -- (2.75,0) -- (4.25,1.2) -- (5.75,0) -- (6.0,0.2);
  \draw[figvecres] (6.0,0.2) -- (6.6,0.68);

  \draw[figgray, line width=0.55pt] (3.5,-0.28) -- (3.5,0.35);
  \node[figlblaux, anchor=north] at (3.5,-0.3) {núcleo: $n_1$};

  \node[figlbl, anchor=north, align=center] at (3.0,-0.9)
    {\textbf{(a)} la luz sigue el cable sin pérdidas,\\
     siempre que no se lo curve demasiado};
\end{scope}

% ---------------- (b) onda evanescente ----------------
\begin{scope}[yshift=-4.6cm]
  \fill[figblue!10] (-2.8,-2.0) rectangle (3.4,0);
  \fill[figblue!10] (-2.8,0.95) rectangle (3.4,2.15);
  \draw[figline, line width=1pt] (-2.8,0) -- (3.4,0);
  \draw[figline, line width=1pt] (-2.8,0.95) -- (3.4,0.95);
  \node[figlblaux, anchor=west] at (-2.75,-0.3) {medio 1: $n_1$};
  \node[figlblaux, anchor=east] at (3.35,0.2) {medio 2: $n_2 < n_1$};
  \node[figlblaux, anchor=west] at (-2.75,1.25) {tercer medio};

  % rayo incidente y reflejado
  \draw[figvecres] (-2.3,-1.55) -- (0,0);
  \draw[figvecres] (0,0) -- (2.3,-1.55);

  % onda evanescente: amplitud que decae exponencialmente hacia arriba
  \draw[figaux] (0,0) -- (0,0.95);
  \draw[figamber, line width=1pt, smooth, samples=40, variable=\yy, domain=0:0.95]
    plot ({1.15*exp(-2.6*\yy)}, {\yy});
  \node[figlbl, figamber, anchor=east, inner sep=2pt] at (-0.15,0.5)
    {amplitud $\propto e^{-y/\delta}$};

  % lo que sigue del otro lado
  \draw[figvecres, figred!55] (0,0.95) -- (1.35,1.85);
  \node[figlbl, figred!70, anchor=west, inner sep=2pt] at (1.45,1.7)
    {onda transmitida (débil)};

  \draw[figgray, line width=0.55pt,
        {Latex[length=1.4mm,width=1mm]}-{Latex[length=1.4mm,width=1mm]}]
    (-1.9,0) -- (-1.9,0.95);
  \node[figlblaux, anchor=east, inner sep=2pt] at (-1.94,0.48) {$e$};

  \node[figlbl, anchor=north, align=center] at (0.3,-2.15)
    {\textbf{(b)} con espesor $e$ finito la onda continúa:\\
     reflexión total \emph{frustrada}};
\end{scope}

\end{tikzpicture}\\[0.5em]
{\small\itshape La fibra óptica resuelve un problema concreto: la señal viaja
mejor en luz que en cobre ---que se calienta y tiene impurezas--- pero la luz, a
priori, va en línea recta. La reflexión total la obliga a doblar: adentro del
núcleo cada rebote ocurre con incidencia \textbf{rasante}, muy por encima del
ángulo crítico, y por eso no se pierde nada en cada uno. De ahí también la
recomendación de no doblar una fibra bruscamente: al curvarla el ángulo de
incidencia empeora, en algún punto baja de $\theta_c$ y la luz se escapa. Las
fibras reales no son vidrio-aire sino dos vidrios, y la capa exterior tiene un
degradé de índice que da margen: si el rayo se pasa un poco, en vez de perderse
se curva y vuelve hacia adentro. El panel (b) corrige la frase
\enquote{no se transmite nada}: del otro lado sí penetra campo, sólo que no en
forma de onda que se propague, sino como una \textbf{onda evanescente} cuya
amplitud cae exponencialmente. Si el segundo medio no llega hasta el infinito
sino que tiene un espesor finito, y se coloca un tercer medio antes de que la
evanescente se apague, la onda \textbf{continúa} del otro lado ---poco, y
decreciendo exponencialmente con el espesor, pero continúa.}
\end{center}
